% BOOK-FILE-ID: f-72e52ec99863
% section: シグマの計算

\section{シグマの計算}\label{節：シグマの計算}
\paragraphbf{本編とのつながり}
和の記号と添字の役割を確認する.\,\sectionref*{節：階差型}で変化を足し戻し,\,\sectionref*{節：総和型}で総和の式をずらして引くときに使う.

$1^2+2^2+3^2+4^2$のような和は,
\[
  \sum_{k=1}^{4}k^2=1^2+2^2+3^2+4^2=30
\]
と短く書ける.\,$k$を$1$から$4$まで一つずつ動かし,そのたびに$k^2$を足すという意味である.
この$\sum$を\bookterm{sigma}{シグマ}という.
実数の\bookterm{sequence}{数列}$\{a_k\}$\footnote{複素数でも構わない.}に対しても,\,$a_1$から$a_n$までの総和を
\[
  \sum_{k=1}^{n} a_k=a_1+a_2+\cdots+ a_n
\]
と表す.
添字$k$は,和の中で順に動く番号である.同じ範囲なら$\sum_{k=1}^{n}a_k=\sum_{j=1}^{n}a_j$と書き換えても値は変わらない.
一方,上端の$n$は足す項数を決める.例えば$S_n=\sum_{k=1}^{n}a_k$なら
\[
  S_{n+1}=S_n+a_{n+1}
\]
である.
足す項がない空の和は$0$とする.例えば$\sum_{k=1}^{n-1}a_k$は,\,$n=1$のとき$0$である.
この時以下の事柄が一般に成り立つ.
ここで$c$は,和を取る添字$k$に依存しない定数である.
\begin{align*}
  &\sum_{k=1}^{n} c=nc\\
  &\sum_{k=1}^{n} ca_k=c\sum_{k=1}^{n} a_k\\
  &\sum_{k=1}^{n} (a_k+b_k)=\sum_{k=1}^{n} a_k+\sum_{k=1}^{n} b_k
\end{align*}
「定数」は,ここでは$k$を動かしても変わらないという意味である.
例えば$\sum_{k=1}^{n}nk=n\sum_{k=1}^{n}k$では,\,$n$を和の外へ出せる.\,$k$そのものは外へ出せない.
また,一般に以下が成り立つ.
\begin{align*}
  &\sum_{k=1}^{n} k=\frac{1}{2}n(n+1)\\
  &\sum_{k=1}^{n} k^2=\frac{1}{6}n(n+1)(2n+1)\\
  &\sum_{k=1}^{n} k^3=\left\{\frac{1}{2}n(n+1)\right\}^2\\
\end{align*}

\bookterm{recurrence}{漸化式}の添字をずらす場合は,式全体の$n$を置き換える.
例えば$a_{n+1}=2a_n+n$から$a_{n+2}=2a_{n+1}+n+1$を得る.
